A filtration records the leading part of a differential, but it does not determine the differential or a comparison map.
For current algebras, three operations must be distinguished: filtering the full state space, taking sections of the collision sheaves, and lifting a graded map.
We construct the filtration and the lifting equations separately.
The obstruction already occurs in a complex with two basis vectors.

All complexes are cohomologically graded over $\mathbb C$, with differential of degree $1$.
Tensor products of vector spaces are algebraic, and $C[k]^n=C^{n+k}$.
A degree-$a$ map $u:C\to T$ has differential
$\partial u=d_Tu-(-1)^a u d_C$.
Thus a chain map has degree zero and satisfies $\partial u=0$.

\subsubsection{A graded map that has no chain lift}

An increasing filtration on a complex $C$ consists of subcomplexes
$F_{p-1}C\subseteq F_pC$, with $F_{-1}C=0$.
It is exhaustive if every element belongs to some $F_pC$.
Its graded piece is $\operatorname{gr}_p C=F_pC/F_{p-1}C$.
A filtered map sends $F_pC$ into $F_pT$.

\begin{example}[The first lifting obstruction]\label{ex:fl19-no-lift}
Let $C^0=T^0=\mathbb Cx$ and $C^1=T^1=\mathbb Cz$.
Give both spaces the filtration
$F_0=\mathbb Cz$, $F_1=\mathbb Cx\oplus\mathbb Cz$.
Set $d_C=0$, while $d_Tx=z$ and $d_Tz=0$.
Both associated graded differentials vanish, so their graded complexes and first cohomology pages agree.
Any degree-zero map inducing the identity on the graded pieces must send $x$ to $x$ and $z$ to $z$.
The chain equation on $x$ would then give $z=0$, a contradiction.

Equivalently, use the decreasing filtration
$G^0=C$, $G^1=\mathbb Cz$, $G^2=0$.
Then $x$ has filtration index zero and $z$ has index one.
The target has a nonzero first page differential $[x]\mapsto[z]$, whereas the source has none.
Putting $x$ in increasing order $N$ and $z$ in order zero moves the same obstruction to order $N$.
Every earlier differential then vanishes.
\end{example}

\subsubsection{The full current state space}

Let $\mathfrak g$ be a finite-dimensional complex Lie algebra with an invariant symmetric form $\kappa$.
Write $x_m=x\otimes t^m$ and define
\[
\widehat L=\mathfrak g[t,t^{-1}]\oplus\mathbb CK,
\qquad [x_m,y_n]=[x,y]_{m+n}+m\delta_{m+n,0}\kappa(x,y)K.
\]
The element $K$ is central.
For Laurent polynomials $f,g,h$, the central form on $xf,yg$ is
\[
\kappa(x,y)\operatorname{Res}(g\,df).
\]
It is alternating because $\operatorname{Res}d(fg)=0$.
Invariance reduces its cyclic Jacobi sum to
\[
\kappa([x,y],z)\operatorname{Res}
\bigl(h\,d(fg)+f\,d(gh)+g\,d(hf)\bigr)=0,
\]
since the differential in parentheses is $2d(fgh)$.

Put $L_-=t^{-1}\mathfrak g[t^{-1}]$ and
$\mathfrak b=\mathfrak g[t]\oplus\mathbb CK$.
For $k\in\mathbb C$, let $\mathbb C_k$ be the one-dimensional $\mathfrak b$-module on which $\mathfrak g[t]$ acts by zero and $K$ acts by $k$.
It is a module because two nonnegative modes have no central bracket.
The vacuum module is the specified quotient
\[
V_k=U(\widehat L)\otimes_{U(\mathfrak b)}\mathbb C_k,
\qquad \mathbf1=1\otimes1.
\]
Here $U(\mathfrak l)$ is the tensor algebra of a Lie algebra $\mathfrak l$ modulo $xy-yx-[x,y]$.

\begin{lemma}[Ordered current states]\label{lem:fl19-state-pbw}
Order a basis of $L_-$.
The nondecreasing words in this basis, applied to $\mathbf1$, form a basis of $V_k$.
The Poincar\'e--Birkhoff--Witt (PBW) filtration $F_pV_k$ by at most $p$ negative current modes has
\[
\operatorname{gr}_F V_k=\operatorname{Sym}(L_-)
\]
as graded vector spaces.
Give $x_{-n}$ energy $n$ for $n\ge1$.
The energy-$h$ space is finite-dimensional and has PBW order at most $h$.
\end{lemma}
\begin{proof}
For any ordered Lie basis, replace $ba$ with $ab+[b,a]$ when $b>a$.
Each term decreases word length, or preserves length and decreases the number of inversions.
Thus every reduction terminates.
Disjoint replacements commute.
For $a>b>c$, the two overlapping reductions have common ordered term $cba$ and remaining terms
\[
[b,c]a+b[a,c]+[a,b]c,
\qquad c[a,b]+[a,c]b+a[b,c].
\]
Their difference reduces to
$-[a,[b,c]]+[b,[a,c]]+[[a,b],c]=0$.
There are no equal-letter reduction rules in the ordinary Lie case.
Induction on length and inversions, also inside any fixed surrounding word, therefore gives a unique normal form.
Normalization kills every defining relation in every context and fixes ordered words.
Those words are consequently a basis of the quotient.

Apply this argument to a basis of $\widehat L$ with all negative modes preceding the basis of $\mathfrak b$.
Multiplication identifies $U(\widehat L)$, as a right $U(\mathfrak b)$-module, with
$U(L_-)\otimes U(\mathfrak b)$.
Tensoring with $\mathbb C_k$ gives the asserted basis of the actual vacuum module.
Setting each bracket correction to zero gives the symmetric associated graded.
At energy $h$, only modes with $n\le h$ can occur, and a word has at most $h$ factors.
There are finitely many such words because each mode space has dimension $\dim\mathfrak g$.
\end{proof}

PBW order is not energy.
For one current $a$, the states $a_{-2}\mathbf1$ and $a_{-1}^2\mathbf1$ both have energy two, but orders one and two.
Writing $d=\dim\mathfrak g$, ordered-state counting gives
\[
\sum_{h\ge0}\dim(V_k)_h q^h=\prod_{n\ge1}(1-q^n)^{-d}.
\]
Every coefficient uses only finitely many factors, so this is an identity of formal series.
In particular, $\dim(V_k)_2=d+\binom{d+1}{2}$, which is $9$ for $\mathfrak{sl}_2$.
For $b$ full state factors, the corresponding series is the $b$th power of this product.
For the positive-energy subspace $V_{k,+}=\bigoplus_{h>0}(V_k)_h$, it is
$\bigl(\prod_{n\ge1}(1-q^n)^{-d}-1\bigr)^b$.
These are counts of state tensors, before any geometric coefficient functions or forms are included.

The same ordered-basis argument applies to the universal Virasoro vacuum module.
Use generators $L_n,C$ with
$[L_m,L_n]=(m-n)L_{m+n}+\delta_{m+n,0}(m^3-m)C/12$.
The noncentral Jacobi identity follows by expanding the three quadratic coefficients.
For the central part, substitute $r=-m-n$ in
$(m-n)((m+n)^3-(m+n))+(n-r)((n+r)^3-(n+r))+(r-m)((r+m)^3-(r+m))$; the result is zero.
Induce from the subalgebra spanned by $C$ and $L_n$ for $n\ge-1$, with $C$ acting by $c$ and those $L_n$ killing the vacuum.
This character respects brackets, including $[L_{-1},L_1]=-2L_0$.
The ordered monomials in $L_{-n}$, $n\ge2$, form its basis.
Their energy series is $\prod_{n\ge2}(1-q^n)^{-1}$, and their order in energy $h$ is at most $\lfloor h/2\rfloor$.

\begin{example}[State counts do not bound collision sections]\label{ex:fl19-poles}
On a punctured normal coordinate $t=z_1-z_2$, the localized coefficient sheaf permits powers $t^{-N}$.
Give $t$ weight one and $d\log t$ weight zero.
For a nonzero current $a$ and $N\ge1$, the full two-state expression
\[
t^{-(N+1)}
 \bigl(a_{-1}^{N}\mathbf1\otimes a_{-1}\mathbf1\bigr)\,d\log t
\]
has total rotation weight zero: the states contribute $N+1$, the coefficient contributes $-N-1$, and $d\log t$ has weight zero.
Its PBW order is $N+1$.
The basis in Lemma~\ref{lem:fl19-state-pbw} proves that these orders are unbounded and their symbols are nonzero.
Thus a fixed total weight need not bound the PBW filtration of a localized collision sheaf.
This example concerns local sections on the separated-point locus, not a claim about global sections on a compactification.
Any finite global bound must also control the permitted poles and extensions.
\end{example}

There is another distinction between states and an augmentation ideal.
An augmentation of a unital current vertex algebra to the trivial vertex algebra $\mathbb C\mathbf1$ must preserve its products and vacuum.
The first current product is $a_{(1)}b=a_1(b_{-1}\mathbf1)=k\kappa(a,b)\mathbf1$, by the mode commutator and $a_1\mathbf1=0$.
If the augmentation kills the currents, preservation of this product implies $0=k\kappa(a,b)$.
At nonzero level and nonzero $\kappa$, such an augmentation does not exist.
The positive-energy vector subspace above is therefore not automatically a chiral augmentation ideal.
A curved or relative bar construction requires its specified coefficient and curvature data.

\subsubsection{Filtering the full collision carrier}

Let $U$ be a coordinate chart with labeled normal-crossing divisors $D_e$, $e\in E$, where $E$ is finite, and choose normal coordinates $x_e$ with $D_e=\{x_e=0\}$.
For $S\subseteq E$, write $U_S=\bigcap_{e\in S}D_e$, $i_S:U_S\to U$, and
$D^S=\bigcup_{e\notin S}(D_e\cap U_S)$.
Suppose $M_S$ are complexes of coefficient vector spaces with maps
$\mu_{S,e}:M_S\to M_{S\cup\{e\}}$ of degree zero.
Assume these maps commute with the differentials and that the two composites for distinct remaining edges agree.
The full residue carrier is
\begin{equation}\label{eq:fl19-full-residue}
\mathcal R(M)=\bigoplus_{S\subseteq E}
i_{S*}\bigl(\underline{M_S}\otimes\Omega^\bullet_{U_S}(\log D^S)\bigr)[-2|S|].
\end{equation}
The underline denotes the constant coefficient sheaf.
For $v$ of degree $q$ and a form $\alpha$ of degree $r$, the total degree is $q+r+2|S|$.
With the left-normal residue $r_e(d\log x_e\wedge\beta+\gamma)=\beta|_{U_{S\cup\{e\}}}$, put
\[
D(v\otimes\alpha)=d_Mv\otimes\alpha
 +(-1)^qv\otimes d\alpha
 +\sum_{e\notin S}(-1)^q\mu_{S,e}v\otimes r_e\alpha.
\]
Every term has degree one.
The form identities are $r_f d=-d r_f$ and $r_fr_e=-r_er_f$.
Together with the coefficient cochain and square identities, they give $D^2=0$.
The internal mixed terms have signs $(-1)^q$ and $(-1)^{q+1}$.
The two different-residue terms have the same coefficient sign and opposite form signs.

Give every $M_S$ an increasing filtration, and assume $d_M$ and all $\mu_{S,e}$ preserve it.
Filter~\eqref{eq:fl19-full-residue} on the coefficient factor only:
\begin{equation}\label{eq:fl19-residue-filtration}
F_p\mathcal R(M)=\bigoplus_{S\subseteq E}
i_{S*}\bigl(\underline{F_pM_S}\otimes\Omega^\bullet_{U_S}(\log D^S)\bigr)[-2|S|].
\end{equation}
The displayed differential preserves this filtration.
Every subset summand and its shift remain present.
In particular, this is a filtration of the full carrier, not of one normal-form image inside it.

For full state tensors, the coefficient filtration is the sum filtration
\begin{equation}\label{eq:fl19-tensor-filtration}
F_p(A_1\otimes\cdots\otimes A_b)
=\sum_{p_1+\cdots+p_b\le p}
\operatorname{im}(F_{p_1}A_1\otimes\cdots\otimes F_{p_b}A_b).
\end{equation}
The index $p$ measures state order, independently of cohomological degree and the number of blocks.
For a graded tensor carrier $K=B\otimes A$, this says
\[
F_pK^n=\sum_{a+b=n}\ \sum_{i+j\le p}
\operatorname{im}(F_iB^a\otimes F_jA^b\to K^n).
\]
This formula applies only after that underlying tensor carrier has been specified.

\begin{lemma}[Symbols of tensor coefficients]\label{lem:fl19-tensor-symbol}
For filtered complex vector spaces with $F_{-1}=0$, the canonical map
\[
\bigoplus_{p_1+\cdots+p_b=p}
\operatorname{gr}_{p_1}A_1\otimes\cdots\otimes\operatorname{gr}_{p_b}A_b
\longrightarrow \operatorname{gr}_p(A_1\otimes\cdots\otimes A_b)
\]
is an isomorphism.
The same assertion holds for tensor products of sheaves over their coefficient ring if their filtrations split locally through each finite order.
If all coefficient maps and differentials are filtered, the isomorphism intertwines their symbols.
\end{lemma}
\begin{proof}
Choose complements $A_{i,r}$ to $F_{r-1}A_i$ in $F_rA_i$.
Up to order $p$, the tensor product is the direct sum of the tensors with $\sum r_i\le p$.
Quotienting by order $p-1$ leaves precisely $\sum r_i=p$.
The resulting map sends a tensor of classes to the class of its tensor, so it is independent of the complements.
Local splittings give the identical argument for sheaves.
Their canonical maps agree on overlaps and therefore glue.
Applying any filtered map before or after passing to classes gives the same element, proving compatibility.
\end{proof}

For chiral coefficients, the actual input includes state sheaves away from collisions and maps to diagonal-supported targets.
These maps need not be multiplication of two vectors in a fixed fiber.
To apply the preceding lemma to a chiral bar model, one must supply the full coefficient sheaves, their localized or extended domains, and their typed collision maps.
Each map must preserve~\eqref{eq:fl19-tensor-filtration}, including every lower-pole and regular contribution selected by its test form.
The coordinate transitions, coefficient connection and twisting terms must preserve the same filtration.
Their symbols must be compared with the proposed graded differential on those same sheaves.
Flatness of the state fibers does not construct any of these maps.

For current pairs, the regular normally ordered product has order two, the bracket coefficient has order one, and the central coefficient has order zero.
Consequently the leading operator in the order filtration is determined by the test-form residue and its degree in state order.
It is not determined by the order of the OPE pole.
A bracket that lowers state order appears in a lower filtration component, rather than in the order-preserving differential.

\subsubsection{The twisting term and its square}

The notation for a twisted tensor product also requires maps.
Let $B$ be a cochain complex with a coassociative coproduct
$\Delta:B\to B\otimes B$ of degree zero, and assume that $\Delta$ commutes with the tensor differential.
Let $A$ be a differential graded associative algebra.
For a degree-one map $\tau:B\to A$, write
\[
\partial\tau=d_A\tau+\tau d_B,
\qquad \tau\star\tau=\mu_A(\tau\otimes\tau)\Delta.
\]
The tensor sign in the second expression is
$(\tau\otimes\tau)(c'\otimes c'')=(-1)^{|c'|}\tau(c')\otimes\tau(c'')$.
Assume $\partial\tau+\tau\star\tau=0$.
On the specified carrier $B\otimes A$, define
\begin{equation}\label{eq:fl19-twisting-differential}
d_\tau(c\otimes a)=d_Bc\otimes a+(-1)^{|c|}c\otimes d_Aa
 +\sum_{\Delta c}(-1)^{|c'|}c'\otimes\tau(c'')a.
\end{equation}
Every coproduct value is a finite sum in the algebraic tensor product.

\begin{lemma}[Filtered twisting]\label{lem:fl19-filtered-twisting}
The operator~\eqref{eq:fl19-twisting-differential} squares to zero.
Suppose the differentials and $\tau$ preserve their filtrations, and
\[
\Delta(F_pB)\subseteq\sum_{i+j\le p}F_iB\otimes F_jB,
\qquad F_iA\,F_jA\subseteq F_{i+j}A.
\]
Then $d_\tau$ preserves the sum filtration on $B\otimes A$.
Its symbol is the twisted differential of the symbols whenever the tensor symbol maps are isomorphisms.
\end{lemma}
\begin{proof}
Let $d_0$ be the untwisted tensor differential and let $T_\tau=d_\tau-d_0$.
The tensor Leibniz signs give
$d_0T_\tau+T_\tau d_0=T_{\partial\tau}$.
Here an operator $T_u$ for a degree-two map $u$ has no sign when moved past $c'$.
Coassociativity gives $T_\tau^2=T_{\tau\star\tau}$:
on a triple coproduct $c_1\otimes c_2\otimes c_3$, the two successive signs multiply to $(-1)^{|c_2|}$.
This is exactly the sign in the product of the two $\tau$ values.
Thus $d_\tau^2=T_{\partial\tau+\tau\star\tau}=0$.
For filtration, a term of orders $i,j$ in $\Delta c$ and order $l$ in $a$ gives output order at most $i+j+l$.
All other terms preserve order by hypothesis.
Taking classes and using Lemma~\ref{lem:fl19-tensor-symbol} proves the final assertion.
\end{proof}

\begin{example}[A twisted polynomial resolution]\label{ex:fl20-twisted-polynomial}
Take $B=\mathbb C1\oplus\mathbb Cc$ with $|c|=-1$, zero differential, and
$\Delta c=1\otimes c+c\otimes1$, $\Delta1=1\otimes1$.
Take $A=\mathbb C[x]$ in degree zero, with zero differential, and set $\tau(c)=x$, $\tau(1)=0$.
The coproduct is coassociative, and both terms in the twisting equation vanish.
Formula~\eqref{eq:fl19-twisting-differential} gives
\[
d_\tau(c\otimes x^r)=1\otimes x^{r+1},
\qquad d_\tau(1\otimes x^r)=0.
\]
Give $c$ and $x$ order one, and $1$ order zero.
The differential preserves total order.
Define $H(1\otimes x^r)=c\otimes x^{r-1}$ for $r\ge1$, and let $H$ vanish on $1\otimes1$ and on $c\otimes A$.
Then $d_\tau H+Hd_\tau$ is the identity minus projection to $\mathbb C(1\otimes1)$.
This is a filtered contraction onto the constant class on the specified algebraic carrier.
\end{example}

This lemma specifies the algebraic signs and filtration requirements for a twisting term.
In a chiral model, the coproduct, coefficient action and twisting map must have their actual sheaf domains and diagonal targets.
The notation $\otimes_\tau$ does not construct those maps or prove their equation.
If the equation has a nonzero defect, the same calculation gives the corresponding operator $d_\tau^2$; the carrier is a cochain complex only if that defect acts by zero.

\subsubsection{Sections and their graded lifting obstruction}

Let $\mathcal K^n$ be a sheaf with filtration $F_p\mathcal K^n$, and use honest global sections
$K^n=\Gamma(Y,\mathcal K^n)$.
Set $F_pK^n=\Gamma(Y,F_p\mathcal K^n)$.
The natural map
\begin{equation}\label{eq:fl19-global-symbol}
\operatorname{gr}_p\Gamma(Y,\mathcal K^n)
\longrightarrow\Gamma(Y,\operatorname{gr}_p\mathcal K^n)
\end{equation}
is injective, but need not be surjective.
For a global symbol $s$, choose local lifts $a_i$ on a cover $U_i$.
The differences $a_j-a_i$ belong to $F_{p-1}\mathcal K^n$ and satisfy the cocycle equation on triple intersections.
Changing the lifts changes this cocycle by a coboundary.
Its class vanishes, possibly after refining the cover, exactly when the lifts can be adjusted to glue.
This gives the exact obstruction to surjectivity of~\eqref{eq:fl19-global-symbol}.
Here a degree-one \v Cech class means such overlap cocycles modulo coboundaries and refinement.

\begin{example}[A locally split filtration with no global lift]\label{ex:fl19-global-sections}
For algebraic sheaves on $\mathbb P^1$ with homogeneous coordinates $[X:Y]$, use
\[
0\longrightarrow\mathcal O(-2)
\xrightarrow{(-Y,X)}\mathcal O(-1)^{\oplus2}
\xrightarrow{(X,Y)}\mathcal O\longrightarrow0.
\]
The two sections $X,Y$ never vanish together, so the sequence is exact and locally split.
Filter the middle bundle by $F_0=\mathcal O(-2)$ and $F_1=\mathcal O(-1)^{\oplus2}$.
Every graded piece is locally free.
Nevertheless, $\Gamma(\mathcal O(-1)^{\oplus2})=0$, while $\Gamma(\mathcal O)=\mathbb C$.

The failure can be computed without a dimension formula for sheaf cohomology.
On $X\ne0$, write $z=Y/X$ and choose the lift $(1,0)$ in the $X^{-1}$ frame.
On $Y\ne0$, choose $(0,1)$ in the $Y^{-1}$ frame, which becomes $(0,z^{-1})$ on the overlap.
Their difference $(-1,z^{-1})$ equals $z^{-1}(-z,1)$.
Thus the obstruction is $z^{-1}$ in the $X^{-2}$ frame of $\mathcal O(-2)$.
Coboundaries lie in
$\mathbb C[z]+z^{-2}\mathbb C[z^{-1}]$, which has no $z^{-1}$ coefficient.
The class is nonzero.
The same transition description shows that a global section of $\mathcal O(-1)$ would lie in
$\mathbb C[z]\cap z^{-1}\mathbb C[z^{-1}]=0$.
\end{example}

Replacing $\Gamma$ by a \v Cech or derived section complex changes the carrier.
Such a change can be useful, but it requires its own comparison with the stated native complex.
Even exhaustive sheaf filtrations need not give exhaustive filtrations on global sections without a uniform finite order for each section.
Both conditions must therefore be checked on the chosen section carrier.

\subsubsection{A filtered criterion with its exact hypotheses}

\begin{theorem}[Filtered acyclicity]\label{thm:fl19-filtered-acyclicity}
Let $(K,d)$ have an exhaustive increasing filtration with $F_{-1}K=0$ and $dF_p\subseteq F_p$.
If $H^n(\operatorname{gr}_pK)=0$ for every $p$, then $H^n(K)=0$.
No finite-dimensionality or completion assumption is needed.
If a filtered chain map $f:C\to T$ has a quasi-isomorphism on associated graded complexes, it is a quasi-isomorphism.
\end{theorem}
\begin{proof}
A cycle $x$ has finite order $p$.
Its symbol is a graded cycle, hence the differential of a symbol of degree $n-1$.
Choose a lift $y\in F_pK^{n-1}$ of that primitive.
Then $x-dy$ has order at most $p-1$.
Repeat until order $-1$, which is zero.
The finite sum of the chosen $y$ is a primitive for $x$.

For the map assertion use
$\operatorname{Cone}(f)^n=T^n\oplus C^{n+1}$ with differential
$D(y,x)=(d_Ty+fx,-d_Cx)$ and the sum filtration.
Its graded complex is the cone of $\operatorname{gr}f$ and is acyclic.
The first assertion makes the full cone acyclic.
For a target cycle $y$, a primitive of $(y,0)$ expresses its class as $f[x]$, proving surjectivity.
If $x$ is a source cycle with $fx=d_Ty$, then $(-y,x)$ is a cone cycle.
A cone primitive makes $x$ a source boundary, proving injectivity.
\end{proof}

\begin{theorem}[PBW criterion on a supplied chiral coefficient complex]
\label{thm:pbw-koszulness-criterion}
Let $K=\bar B^{\mathrm{ch}}(\mathcal A)\otimes_\tau\mathcal A$ denote a specified native cochain complex, including its full state sheaves, section functor, collision maps and twisting differential.
Assume the following conditions.
\begin{enumerate}
\item\label{item:pbw-flat} Its state-order filtration is~\eqref{eq:fl19-tensor-filtration} on every coefficient tensor and includes all collision strata with their grading shifts.
The actual differential and transitions preserve it.
The tensor symbol maps are isomorphisms on the supplied coefficient sheaves.
\item\label{item:pbw-classical-koszul} After the prescribed section functor, the graded symbol complex has a specified, order-preserving chain isomorphism or quasi-isomorphism to a graded complex $K_0$ whose positive-degree cohomology vanishes in every order.
This is a comparison of full coefficient complexes and their operations, not only of state vector spaces.
\item\label{item:pbw-bounded} On the chosen global section carrier, the filtration is exhaustive, $F_{-1}K=0$, and the global symbol maps~\eqref{eq:fl19-global-symbol} are isomorphisms.
\end{enumerate}
Then $H^n(K)=0$ for $n>0$.
Thus this specified native Koszul complex is acyclic in the degrees required for chiral Koszulness.
\end{theorem}
\begin{proof}
The first condition defines a filtered cochain complex on the stated carrier.
The third condition identifies its graded section complex with the sections of the symbol complex.
In the second condition, the comparison and vanishing are understood on this same section carrier, so they give
$H^n(\operatorname{gr}_pK)=0$ for every $p$ and $n>0$.
Theorem~\ref{thm:fl19-filtered-acyclicity} applies.
In particular, no finite-dimensional state count is substituted for a global section comparison or a filtration bound.
\end{proof}

The classical polynomial resolution is one possible graded complex, but it must be identified by an actual map.
For a vector space $V$ with basis $x_i$, put
$P=\operatorname{Sym}(V)\otimes\Lambda(e_i)$, with $|x_i|=0$, $|e_i|=-1$ and $d e_i=x_i$.
The differential extends by the graded Leibniz rule and squares to zero.
The degree-minus-one derivation $h$ defined by $h(x_i)=e_i$ and $h(e_i)=0$ satisfies
$dh+hd=N$ on monomials of total polynomial and exterior length $N$.
Indeed its anticommutator is a degree-zero derivation fixing every generator.
Consequently $h/N$ contracts every positive length, leaving only the constant class.
Each expression is finite even for an infinite basis.
This proves polynomial Koszul exactness internally, but supplies no map from a localized chiral coefficient complex to $P$.

\begin{example}[Completion changes the acyclicity question]\label{ex:fl19-completion}
Take $C^0=C^1=\mathbb C[u]$, differential multiplication by $1-u$, and decreasing filtration $G^p=u^p\mathbb C[u]$.
Every graded differential is the identity, so the graded complex is acyclic.
The polynomial complex nevertheless has cokernel $\mathbb C$, by evaluation at $u=1$.
In the completion $\mathbb C[[u]]$, the differential is invertible with inverse $\sum_{j\ge0}u^j$.
The polynomial filtration is separated but not complete.
The increasing polynomial-degree filtration does not preserve this differential, so Theorem~\ref{thm:fl19-filtered-acyclicity} does not apply to it.
\end{example}

\subsubsection{The chain equation at each filtration order}

Let $C,T$ be cochain complexes with exhaustive increasing filtrations starting at order zero.
Choose degree-preserving splittings
\[
C=\bigoplus_{p\ge0}C_p,\qquad T=\bigoplus_{p\ge0}T_p,
\qquad F_pC=\bigoplus_{i\le p}C_i,
\]
and similarly for $T$.
The differentials have unique components
\[
d_C=\sum_{r\ge0}d_{C,r},\qquad d_{C,r}:C_p^n\to C_{p-r}^{n+1},
\]
and likewise for $T$.
All terms with negative target order are zero, so these sums are finite on each input.
The identities $d_C^2=d_T^2=0$ give
$\sum_{i+j=r}d_{C,i}d_{C,j}=0$ and the corresponding target identities.

For $r\ge0$, let
\[
\mathcal H_r^a=\prod_{p,n}\operatorname{Hom}_{\mathbb C}(C_p^n,T_{p-r}^{n+a}),
\qquad \partial_0u=d_{T,0}u-(-1)^a u d_{C,0}.
\]
This is a cochain complex, since the two order-zero differentials square to zero.
Fix a graded chain map $f_0\in\mathcal H_0^0$.
Suppose $f_1,\ldots,f_{r-1}$ have already solved the chain equation in orders below $r$.
The next defect is the explicit degree-one map
\begin{equation}\label{eq:fl19-obstruction}
\omega_r=\sum_{i=1}^{r}
\bigl(d_{T,i}f_{r-i}-f_{r-i}d_{C,i}\bigr)\in\mathcal H_r^1.
\end{equation}

\begin{theorem}[The exact lifting obstruction]\label{thm:fl19-lift-obstruction}
The map $\omega_r$ is $\partial_0$-closed.
The chosen partial lift extends through order $r$ exactly when
$[\omega_r]=0$ in $H^1(\mathcal H_r,\partial_0)$.
An extension is a map $f_r\in\mathcal H_r^0$ satisfying
\begin{equation}\label{eq:fl19-lift-equation}
\partial_0 f_r=-\omega_r.
\end{equation}
All extensions differ by degree-zero cycles in $\mathcal H_r$.
If compatible choices can be made at every order, the sum $F=\sum_{r\ge0}f_r$ is an actual filtered chain map $C\to T$ inducing $f_0$.
\end{theorem}
\begin{proof}
Put $F_{<r}=\sum_{j<r}f_j$ and $\Omega=d_TF_{<r}-F_{<r}d_C$.
The square-zero identities give $d_T\Omega+\Omega d_C=0$.
The components of $\Omega$ below order $r$ vanish by the induction hypothesis.
Its order-$r$ component is~\eqref{eq:fl19-obstruction}, and the first component of the displayed identity is $\partial_0\omega_r=0$.
Adding $f_r$ changes that defect to $\omega_r+\partial_0f_r$.
This proves the obstruction assertion and the description of all extensions.

For an input in $F_pC$, all $f_r$ with $r>p$ vanish.
Consequently $F$ is defined by a finite sum on every input.
Its chain defect vanishes in every order, so it is zero.
No completion or convergence of an infinite series on a single input is used.
\end{proof}

The first obstruction in Example~\ref{ex:fl19-no-lift} is the map $x\mapsto z$ in $\mathcal H_1^1$.
There $\partial_0=0$, so it is not a boundary.
Higher obstruction classes depend on the preceding choices of lifts.
Their vanishing must hold along one compatible sequence of choices, not merely for unrelated partial lifts.

We next give actual primitives when the obstruction vanishes.
For each order-zero complex, put $B^n=\operatorname{im}d_0^{n-1}$ and $Z^n=\ker d_0^n$.
Choose splittings
$Z^n=B^n\oplus H^n$ and $C^n=Z^n\oplus A^n$ within each order.
The map $d_0:A^n\to B^{n+1}$ is an isomorphism.
Let $h_C$ be its inverse on boundaries and zero on $H^n$ and $A^n$.
Let $\pi_C$ project to $H$, and $\iota_C$ include $H$ into $C$.
Then
\[
d_{C,0}h_C+h_Cd_{C,0}=1-\iota_C\pi_C.
\]
This follows separately on boundaries, the chosen cohomology representatives, and $A$.
Make the same choices for $T$.
These are splittings of vector spaces, with no implied linearity over a sheaf of functions or differential operators.

For a homogeneous $u\in\mathcal H_r^a$, define
\begin{equation}\label{eq:fl19-hom-contraction}
S(u)=h_Tu+(-1)^a\iota_T\pi_Tu h_C,
\qquad \Pi(u)=\pi_Tu\iota_C.
\end{equation}
The first map has degree $-1$ on the Hom complex, while the second takes values in maps between the chosen graded cohomology spaces.

\begin{lemma}[A primitive for a vanishing defect]\label{lem:fl19-hom-primitive}
The identities
\[
\partial_0 S+S\partial_0=1-\bigl(u\mapsto\iota_T\Pi(u)\pi_C\bigr),
\qquad \Pi\partial_0=0
\]
hold on every $\mathcal H_r$.
A closed $\omega\in\mathcal H_r^1$ is a boundary if and only if $\Pi(\omega)=0$.
When this holds, $-S(\omega)$ solves~\eqref{eq:fl19-lift-equation}.
\end{lemma}
\begin{proof}
The anticommutator with $u\mapsto h_Tu$ is $(1-\iota_T\pi_T)u$.
For the second term in~\eqref{eq:fl19-hom-contraction}, the signs $(-1)^a$ and the sign in the Hom differential give
\[
\iota_T\pi_Tu(h_Cd_{C,0}+d_{C,0}h_C)
=\iota_T\pi_Tu(1-\iota_C\pi_C).
\]
Adding the two expressions proves the first identity.
Projection kills target boundaries, and the source inclusion consists of cycles, proving $\Pi\partial_0=0$.
For a closed $\omega$ with zero projection, the first identity gives $\omega=\partial_0S(\omega)$.
The converse follows from the second identity.
\end{proof}

\begin{example}[A nonzero correction that gives a lift]\label{ex:fl19-successful-lift}
Let $C^0=T^0=\mathbb Cx\oplus\mathbb Ca$ and $C^1=T^1=\mathbb Cz$.
Give $a,z$ order zero and $x$ order one.
Set $d_Ca=d_Ta=z$, $d_Cx=0$ and $d_Tx=z$.
The identity is a chain map for the order-zero differentials.
Its first defect sends $x$ to $z$, which is now a boundary in the target graded complex.
Choose $h_Tz=a$ and let $h_T$ vanish on the other basis vectors.
The correction is $f_1x=-a$, with $f_1a=f_1z=0$.
Consequently
\[
F(x)=x-a,\qquad F(a)=a,\qquad F(z)=z
\]
is a filtered chain map: $d_T(x-a)=0$.
Its inverse sends $x$ to $x+a$ and fixes $a,z$.
The extra primitive $a$ is precisely the datum absent from Example~\ref{ex:fl19-no-lift}.
\end{example}

\begin{corollary}[An explicit sufficient condition for lifting]\label{cor:fl19-explicit-lift}
Suppose
\[
\operatorname{Hom}\bigl(H^n(C_p,d_{C,0}),
H^{n+1}(T_{p-r},d_{T,0})\bigr)=0
\quad\text{for all }p,n\text{ and }r\ge1.
\]
Then the recursion $f_r=-S(\omega_r)$ constructs a filtered chain lift of every graded chain map $f_0$.
If $f_0$ is a graded quasi-isomorphism, the resulting $F$ is a quasi-isomorphism.
\end{corollary}
\begin{proof}
The displayed condition makes the degree-one cohomology projection $\Pi(\omega_r)$ zero for every closed defect.
The preceding lemma therefore supplies the correction at every order.
Theorem~\ref{thm:fl19-lift-obstruction} constructs the chain map, and
Theorem~\ref{thm:fl19-filtered-acyclicity} proves the quasi-isomorphism assertion.
\end{proof}

For example, the sufficient condition holds when every graded cohomology space of both complexes is concentrated in degree zero.
It does not follow from an isomorphism between their cohomology spaces in several degrees.
An identification of those spaces can first be made into a graded vector-space chain map by
$f_0=\iota_T f_H\pi_C$, where $f_H$ is the proposed cohomology map.
The positive-order obstructions still have to vanish.

\subsubsection{Explicit homotopies and complete filtrations}

The quasi-isomorphism conclusion has a constructive form.
Let $F:C\to T$ be a filtered chain lift of a graded quasi-isomorphism, and let $(M,D)=\operatorname{Cone}(F)$.
Choose an order-preserving contraction $h_0$ of its acyclic graded complex, using the boundary splittings above.
Through the chosen splitting of $M$, write $D=D_0+\delta$.
The operator $\delta$ strictly lowers order.
Put
\begin{equation}\label{eq:fl19-cone-contraction}
N=-\delta h_0-h_0\delta,
\qquad H=h_0\sum_{j\ge0}N^j.
\end{equation}
Every sum is finite on each input, because $N$ lowers order and $F_{-1}M=0$.
Since $Dh_0+h_0D=1-N$ and $D^2=0$, the operator $N$ commutes with $D$.
It follows that
\[
DH+HD=(1-N)\sum_{j\ge0}N^j=1.
\]
Thus~\eqref{eq:fl19-cone-contraction} is an actual filtered contraction of the cone.

Write its blocks on $T\oplus C[1]$ as
$H=\left(\begin{smallmatrix}A&B\\G&E\end{smallmatrix}\right)$.
The equation $DH+HD=1$ gives
\[
d_CG=Gd_T,\qquad
FG=1-d_TA-Ad_T,\qquad
GF=1+d_CE+Ed_C.
\]
Consequently $G$ is a filtered chain inverse up to homotopy.
The homotopies for $1-FG$ and $1-GF$ are $A$ and $-E$, respectively.
No coalgebra property or geometric naturality is inferred for these chosen vector-space maps.

There is a separate complete version.
Suppose instead that the chosen carriers are products
$C=\prod_{p\ge0}C_p$ and $T=\prod_{p\ge0}T_p$, with decreasing filtrations
$G^pC=\prod_{j\ge p}C_j$.
Assume every positive-order differential component raises $p$.
Define the order-$r$ Hom maps to raise $p$ by $r$ as well.
The obstruction equations and their proofs are unchanged.
The sums of corrections and the geometric series in~\eqref{eq:fl19-cone-contraction} converge componentwise: an output of order $p$ receives only orders $0,\ldots,p$.
They define continuous maps of these specified complete, separated products.
The same contraction equations hold in each quotient by $G^p$, hence on the product.
This proves quasi-isomorphism without interchanging cohomology with an unspecified inverse limit.
It makes no assertion about an uncompleted section carrier such as Example~\ref{ex:fl19-completion}.

\subsubsection{Which semi-infinite target is being compared}

The target also requires actual operators.
Let $\mathfrak l=\bigoplus_{m\in\mathbb Z}\mathfrak l_m$ be an ordinary graded Lie algebra with finite-dimensional components and chosen homogeneous bases $e_{a,m}$.
Introduce odd operators $c^{a,m}$ of cohomological degree $1$ and $b_{a,m}$ of degree $-1$, with
\[
\{b_{a,m},c^{b,n}\}=\delta_a^b\delta_m^n,
\qquad \{b,b\}=\{c,c\}=0.
\]
The braces mean $uv+vu$ for odd operators.
Fix a vacuum killed by $c^{a,m}$ for $m\le0$ and by $b_{a,m}$ for $m>0$.
The ghost space consists of finite products of the remaining creation operators on this vacuum.
Equivalently, its underlying vector space is
\[
\mathcal F=\Lambda\bigl(c^{a,m}\ (m>0),\ b_{a,m}\ (m\le0)\bigr).
\]
Creation acts by exterior multiplication, and the paired annihilator acts by exterior contraction with the Koszul sign.
Its degree is measured relative to the vacuum, whose degree is zero.
Assign energy $m$ to $c^{a,m}$ and $-m$ to $b_{a,m}$.
All creation energies are nonnegative, with only finitely many zero-energy generators.
Each fixed-energy ghost space is finite-dimensional.

Let $M=\bigoplus_{h\ge0}M_h$ be an ordinary module, placed in cohomological degree zero, with $\rho(e_{a,m})M_h\subseteq M_{h-m}$.
Set $M_h=0$ for $h<0$ and use the algebraic carrier $M\otimes\mathcal F$.
The linear part of the proposed differential is
$\sum_{a,m}\rho(e_{a,m})c^{a,m}$.
The Lie bracket part has the form
\[
-\tfrac12\sum_{i,j,k} f_{ij}^{k}{:}c^i c^j b_k{:},
\]
where $[e_i,e_j]=\sum_kf_{ij}^ke_k$ and normal ordering moves creation operators to the left with the fermionic signs.
An energy-preserving scalar linear charge term is a specified finite combination of the $c^{a,0}$.

\begin{lemma}[The algebraic operator domain]\label{lem:fl20-ghost-domain}
The linear current term and the normal-ordered cubic term define locally finite degree-one, energy-zero operators on $M\otimes\mathcal F$.
The same holds after adding the stated charge term.
\end{lemma}
\begin{proof}
For $v\in M_h$, the modes with $m>h$ annihilate $v$.
For $m\le0$, the ghost $c^{a,m}$ annihilates a fixed exterior monomial unless that monomial contains $b_{a,m}$.
Thus the linear sum has finitely many nonzero terms on every pure tensor.

In a normal-ordered cubic monomial, all annihilators act before all creators.
Every annihilator must match one of the finitely many occupied creation variables of the input.
Their indices therefore range over a finite set.
The Lie grading gives $m_k=m_i+m_j$, so each cubic term has total energy zero.
Its creation energies are nonnegative, and their sum equals the energy removed by the annihilators.
Only finitely many creation variables have energy at most that fixed sum, including the finite-dimensional zero-mode space.
If a term has no annihilators, its creators all have energy zero, giving the same finiteness conclusion.
This proves local finiteness on each exterior monomial and hence on every algebraic input.
The operator degrees follow from $|c|=1$ and $|b|=-1$.
\end{proof}

The square-zero identity is a separate condition on these actual operators.
It depends on the coefficient action, central generators and polarization.
On a completion allowing infinitely many excitations, the preceding algebraic proof does not establish continuity or convergence.

\begin{example}[The central ghost changes the complex]\label{ex:fl19-central-ghost}
For one Heisenberg current with $[a_m,a_n]=m\delta_{m+n,0}K$, use the polynomial coefficient module $M=\mathbb C[x_1,x_2,\ldots]$ with mode action
\[
a_{-m}=x_m,\qquad a_m=km\partial_{x_m}\quad(m>0),\qquad a_0=0.
\]
Here $K$ acts by $k$ and $x_m$ has energy $m$.
Use ghosts $c^m,b_m$ with the preceding polarization and put
\[
Q_0=\sum_{m\in\mathbb Z}a_m c^m,
\qquad J=\sum_{m>0}m c^m c^{-m}.
\]
Both sums are finite on every polynomial with finitely many ghost excitations.
For positive $m$, $a_m$ kills a fixed polynomial for sufficiently large $m$.
For negative $m$, $c^m$ kills a fixed ghost vector except at its finitely many holes.
The same observation proves local finiteness of $J$.
The commutators of the $a_m$ give
\[
Q_0^2=kJ,
\qquad Q_0^2(1\otimes b_{-1}|0\rangle)
=k\otimes c^1|0\rangle.
\]
The latter vector is nonzero when $k\ne0$.
A scalar linear term in the $c^m$ anticommutes with $Q_0$ and squares to zero, so it does not remove this defect.

Now retain a central Lie generator acting by $k$, and add its ghosts $c^K,b_K$ with $\{c^K,b_K\}=1$.
Give both central ghosts energy zero, and choose $c^K|0\rangle=0$.
Thus the enlarged carrier is $M\otimes\Lambda(c^m\ (m>0),b_m\ (m\le0),b_K)$, with $c^K$ contracting the last exterior variable.
The absolute differential is
\[
Q=Q_0+k c^K-Jb_K.
\]
The new terms square to $-kJ$, and their anticommutator with $Q_0$ vanishes, so $Q^2=0$.
However, $Qb_K+b_KQ=k$.
For $k\ne0$, the map $b_K/k$ contracts the entire complex.
Thus omitting the central ghost gives a nonzero square in this model, while retaining it gives an acyclic absolute complex.
A relative construction or an additional anomaly-cancelling sector is different data.
\end{example}

The example does not identify a general chiral bar complex with either target.
It shows why a proposed semi-infinite comparison must specify its central treatment and differential before computing a page.
Likewise, conformal energy does not identify the full ghost space with exterior powers of the zero modes.
Positive-energy creation operators remain in its graded pieces.
Discarding them requires a proved contraction or another comparison map.

\subsubsection{The native comparison and its descent conditions}

Suppose the full native bar complex $C$ and a specified square-zero semi-infinite target $T$ have been constructed.
Their filtrations must be defined on all state, form and ghost factors.
A proposed initial map must be a chain map
$f_0:\operatorname{gr}C\to\operatorname{gr}T$ on those actual carriers.
An isomorphism to the cohomology of a finite-dimensional Lie algebra is not such a construction.
Given $f_0$, equations~\eqref{eq:fl19-obstruction}--\eqref{eq:fl19-lift-equation} state exactly what each higher collision or mode correction must solve.
If their obstructions vanish compatibly, they give the actual map $F$.
If $f_0$ is a quasi-isomorphism and the stated algebraic or complete filtration hypotheses hold, the cone construction supplies its inverse and homotopies.
These are conditional conclusions about the specified native objects.

If a sheaf or differential-operator comparison is required, the corrections must belong to the corresponding complex of morphisms.
The vector-space splittings in~\eqref{eq:fl19-hom-contraction} need not preserve coefficient multiplication, differential operators or restriction maps.
One may use that formula in such a category only when the displayed projections and homotopies have those compatibilities.
Otherwise the obstruction must be solved in that smaller morphism complex.

For local maps $F_i$ on a cover, a first overlap homotopy has degree $-1$ and must satisfy
\[
\partial H_{ij}=F_j-F_i.
\]
The triple-overlap difference $H_{jk}-H_{ik}+H_{ij}$ is closed.
For the total differential $\delta+(-1)^r\partial$ in \v Cech degree $r$, the next required map has degree $-2$ and satisfies
\[
\partial K_{ijk}=-(H_{jk}-H_{ik}+H_{ij}).
\]
Closedness follows by applying $\partial$ and cancelling the three differences of $F_i$.
Existence of $K$ requires that closed map to be a boundary.
Further overlaps give the analogous equations in lower internal degrees.
Strict descent requires the maps themselves to agree on overlaps; homotopy-coherent descent instead uses these additional maps and a specified total section complex.
Neither replaces the underived global sections silently.

Thus the remaining native construction consists of filtered collision and twisting operators on the full state sheaves, the global symbol comparison, and a specified polarized semi-infinite differential.
One must then provide $f_0$ and solve its filtered and overlap equations in the required coefficient category.
The state-level PBW basis, the polynomial Koszul contraction, and an identification of first pages do not supply these maps.
